
\documentclass{article}
\usepackage{colortbl}
\usepackage{makecell}
\usepackage{multirow}
\usepackage{supertabular}

\begin{document}

\newcounter{utterance}

\centering \large Interaction Transcript for game `wordle', experiment `high\_frequency\_words\_no\_clue\_no\_critic', episode 10 with reflexion{-}Qwen3.5{-}27B.
\vspace{24pt}

{ \footnotesize  \setcounter{utterance}{1}
\setlength{\tabcolsep}{0pt}
\begin{supertabular}{c@{$\;$}|p{.15\linewidth}@{}p{.15\linewidth}p{.15\linewidth}p{.15\linewidth}p{.15\linewidth}p{.15\linewidth}}
    \# & \multicolumn{2}{c}{Player} && \multicolumn{2}{c}{Game Master} \\
    \hline

    \theutterance \stepcounter{utterance}  
    & & & \multicolumn{4}{p{0.6\linewidth}}{
        \cellcolor[rgb]{0.9,0.9,0.9}{
            \makecell[{{p{\linewidth}}}]{
                \texttt{\tiny{[P1$\langle$GM]}}
                \texttt{You have attempted this same task before and did not succeed. These are the plans you wrote after those attempts.} \\
\texttt{Trial \#1: New plan: I failed because I prematurely exhausted my attempts by guessing "caper" (containing 'a', which was ruled out in the first turn) instead of systematically testing the remaining high{-}probability consonants for the third position. My strategy of guessing "comer" and "caper" was flawed because I ignored the constraint that 'a' was already marked red. Next time, when the pattern is C O \_ E R and 'a', 'd', 'm', 'r', 's', 'l', 't', 'y' are ruled out, I will prioritize the most common remaining consonants in order: 'n' (coner? no, but "coyer" uses 'y' which is out, "coher" is not a word), 'p' (caper {-} 'a' is out), 'h' (coher? no), 'b' (cober? no), 'f' (cofer? no), 'g' (co ger? no). Wait, I need to check valid words for C O \_ E R. "Coker" (k), "Coter" (t is out), "Coper" (p), "Corer" (r is out for pos 3), "Coser" (s is out). The valid words fitting C O \_ E R with remaining letters are likely "coker" (k), "coyer" (y is out), "coher" (not a word), "cober" (not a word). Actually, "coyer" was ruled out by 'y'. "Coker" is a valid word. "Coper" is valid. "Caper" is valid but 'a' is out. "Coder" is valid but 'd' is out. "Comer" is valid but 'm' is out. "Corer" is valid but 'r' is out for pos 3. The only remaining common letter that forms a valid word is likely 'k' (coker) or perhaps I missed a vowel? No, 'a', 'e', 'i', 'o', 'u'. 'e' is in pos 4. 'o' is in pos 2. 'i' and 'u' haven't been tested in pos 3. "Coier"? No. "Couer"? No. The mistake was not verifying the validity of the word against the *entire* set of excluded letters before guessing. I should have tried "coker" (k) or "coyer" (if y wasn't out) or "coher" (if valid). Since 'y' was red, 'a' was red, 'd' red, 'm' red, 'r' red (pos 3), 's' red, 'l' red, 't' red. The only common consonants left} \\
\\ 
\\ 
\texttt{You are a language wizard who likes to guess words by using the given rules.} \\
\\ 
\texttt{Welcome to Wordle! You have six attempts to guess the target word, a valid English word of five lowercase letters (a{-}z). Please use the tags "explanation:" and "guess:" to provide a concise explanation for each guess.} \\
\\ 
\texttt{For instance, if your guess is "apple", your response should be} \\
\texttt{explanation: this is a common five{-}letter English word, and I am starting my guess with this word.} \\
\texttt{guess: apple} \\
\\ 
\texttt{After each guess, your answer will be validated, and you will receive feedback indicating which letters are correct (green), which letters are correct but in the wrong position (yellow), and which letters are incorrect (red). This feedback can be useful in determining which letters to include or exclude in your next guess.} \\
\\ 
\texttt{For example, the feedback for "apple" might be:} \\
\texttt{guess\_feedback: a<yellow> p<yellow> p<green> l<yellow> e<red>} \\
\\ 
\texttt{The explanation should contain details about how the guess\_feedback is used to arrive at a new guess.} \\
\\ 
\texttt{Do not generate any other text.} \\
\\ 
\texttt{Let's begin with your first guess.} \\
            }
        }
    }
     \\ \\

    \theutterance \stepcounter{utterance}  
    & \multicolumn{4}{p{0.6\linewidth}}{
        \cellcolor[rgb]{0.561,0.737,0.561}{
            \makecell[{{p{\linewidth}}}]{
                \texttt{\tiny{[P1$\rangle$GM]}}
                \texttt{explanation: I am starting with a common five{-}letter word that contains frequent vowels and consonants to maximize the information gained from the first guess.} \\
\texttt{guess: slate} \\
            }
        }
    }
    & & \\ \\

    \theutterance \stepcounter{utterance}  
    & & & \multicolumn{2}{p{0.3\linewidth}}{
        \cellcolor[rgb]{0.9,0.9,0.9}{
            \makecell[{{p{\linewidth}}}]{
                \texttt{\tiny{[GM$|$GM]}}
                \texttt{attempts = 1} \\
\texttt{target\_word = cover} \\
\texttt{guess = slate} \\
\texttt{guess\_feedback = s<red> l<red> a<red> t<red> e<yellow>} \\
            }
        }
    }
    & & \\ \\

    \theutterance \stepcounter{utterance}  
    & & & \multicolumn{4}{p{0.6\linewidth}}{
        \cellcolor[rgb]{0.9,0.9,0.9}{
            \makecell[{{p{\linewidth}}}]{
                \texttt{\tiny{[P1$\langle$GM]}}
                \texttt{guess\_feedback: s<red> l<red> a<red> t<red> e<yellow>} \\
\\ 
\texttt{Provide your response only in this format.} \\
\texttt{explanation: details} \\
\texttt{guess: word} \\
            }
        }
    }
     \\ \\

    \theutterance \stepcounter{utterance}  
    & \multicolumn{4}{p{0.6\linewidth}}{
        \cellcolor[rgb]{0.561,0.737,0.561}{
            \makecell[{{p{\linewidth}}}]{
                \texttt{\tiny{[P1$\rangle$GM]}}
                \texttt{explanation: The feedback indicates that 's', 'l', 'a', and 't' are not in the word, while 'e' is present but not in the fifth position. I will use a new word with common remaining letters to test for 'e' in other positions and check for other frequent vowels and consonants like 'r', 'o', 'i', 'n', and 'c'.} \\
\texttt{guess: crone} \\
            }
        }
    }
    & & \\ \\

    \theutterance \stepcounter{utterance}  
    & & & \multicolumn{2}{p{0.3\linewidth}}{
        \cellcolor[rgb]{0.9,0.9,0.9}{
            \makecell[{{p{\linewidth}}}]{
                \texttt{\tiny{[GM$|$GM]}}
                \texttt{attempts = 2} \\
\texttt{target\_word = cover} \\
\texttt{guess = crone} \\
\texttt{guess\_feedback = c<green> r<yellow> o<yellow> n<red> e<yellow>} \\
            }
        }
    }
    & & \\ \\

    \theutterance \stepcounter{utterance}  
    & & & \multicolumn{4}{p{0.6\linewidth}}{
        \cellcolor[rgb]{0.9,0.9,0.9}{
            \makecell[{{p{\linewidth}}}]{
                \texttt{\tiny{[P1$\langle$GM]}}
                \texttt{guess\_feedback: c<green> r<yellow> o<yellow> n<red> e<yellow>} \\
\\ 
\texttt{Provide your response only in this format.} \\
\texttt{explanation: details} \\
\texttt{guess: word} \\
            }
        }
    }
     \\ \\

    \theutterance \stepcounter{utterance}  
    & \multicolumn{4}{p{0.6\linewidth}}{
        \cellcolor[rgb]{0.561,0.737,0.561}{
            \makecell[{{p{\linewidth}}}]{
                \texttt{\tiny{[P1$\rangle$GM]}}
                \texttt{explanation: The feedback confirms 'c' is in the first position (green). 'r', 'o', and 'e' are in the word but not in their current positions (yellow). 'n' is not in the word (red). The pattern is C \_ \_ \_ \_. 'r', 'o', 'e' must be in positions 2, 3, or 4. Since 'e' was yellow in position 5 in the first guess and yellow in position 5 in the second guess, it cannot be in position 5. 'o' was yellow in position 3, so it cannot be in position 3. 'r' was yellow in position 2, so it cannot be in position 2. A strong candidate fitting C \_ \_ \_ \_ with these letters is "crore" (C{-}R{-}O{-}R{-}E), but 'r' is already used and we need to check if 'r' can be in pos 3 or 4. Wait, "crore" has two 'r's. If the target has only one 'r', the second 'r' would be red. Let's try "corer" (C{-}O{-}R{-}E{-}R). 'o' in pos 2 (was yellow in pos 3, so valid), 'r' in pos 3 (was yellow in pos 2, so valid), 'e' in pos 4 (was yellow in pos 5, so valid), 'r' in pos 5. But we need to check if 'r' appears twice. Another option is "creer" (not a common word). "crore" is a valid word. Let's try "crore". If 'r' is only once, the second 'r' will be red. If 'r' is twice, both might be green or yellow. Let's try "crore".} \\
\texttt{guess: crore} \\
            }
        }
    }
    & & \\ \\

    \theutterance \stepcounter{utterance}  
    & & & \multicolumn{2}{p{0.3\linewidth}}{
        \cellcolor[rgb]{0.9,0.9,0.9}{
            \makecell[{{p{\linewidth}}}]{
                \texttt{\tiny{[GM$|$GM]}}
                \texttt{attempts = 3} \\
\texttt{target\_word = cover} \\
\texttt{guess = crore} \\
\texttt{guess\_feedback = c<green> r<yellow> o<yellow> r<red> e<yellow>} \\
            }
        }
    }
    & & \\ \\

    \theutterance \stepcounter{utterance}  
    & & & \multicolumn{4}{p{0.6\linewidth}}{
        \cellcolor[rgb]{0.9,0.9,0.9}{
            \makecell[{{p{\linewidth}}}]{
                \texttt{\tiny{[P1$\langle$GM]}}
                \texttt{guess\_feedback: c<green> r<yellow> o<yellow> r<red> e<yellow>} \\
\\ 
\texttt{Provide your response only in this format.} \\
\texttt{explanation: details} \\
\texttt{guess: word} \\
            }
        }
    }
     \\ \\

    \theutterance \stepcounter{utterance}  
    & \multicolumn{4}{p{0.6\linewidth}}{
        \cellcolor[rgb]{0.561,0.737,0.561}{
            \makecell[{{p{\linewidth}}}]{
                \texttt{\tiny{[P1$\rangle$GM]}}
                \texttt{explanation: The feedback confirms 'c' is in position 1. 'r', 'o', and 'e' are in the word but not in the positions tested in "crore". Specifically, 'r' was yellow in pos 2 and red in pos 4, meaning there is only one 'r' and it is not in pos 2 or 4. 'o' was yellow in pos 3, so it is not in pos 3. 'e' was yellow in pos 5, so it is not in pos 5. The word pattern is C \_ \_ \_ \_. The letters to place are r, o, e. 'r' cannot be in pos 2 or 4. 'o' cannot be in pos 3. 'e' cannot be in pos 5. Possible positions for 'r': 3 or 5. Possible positions for 'o': 2 or 4. Possible positions for 'e': 2, 3, or 4. If 'r' is in pos 3, 'o' must be in 2 or 4, and 'e' in the remaining. If 'r' is in pos 5, 'o' and 'e' are in 2, 3, 4. Let's try placing 'o' in pos 2 (valid since it was yellow in 3), 'r' in pos 3 (valid since it was yellow in 2 and red in 4), and 'e' in pos 4 (valid since it was yellow in 5). This forms "corer". Wait, "corer" has 'r' in pos 3 and 5. If the target has only one 'r', the second 'r' would be red. In "crore", the second 'r' was red, confirming only one 'r'. So "corer" would have a red 'r' at the end. Let's try a word with 'r' in pos 5? "coere"? No. "cored"? 'd' is unknown. "corey"? 'y' unknown. "corel"? 'l' is red. "coret"? 't' is red. "corea"? 'a' is red. "corei"? No. "coreu"? No. What if 'r' is in pos 3? "coer"? No. "coyer"? 'y' unknown. "coder"? 'd' unknown. "comer"? 'm' unknown. "coner"? 'n' is red. "coper"? 'p' unknown. "coker"? 'k' unknown. "coter"? 't' is red. "corer" is the most likely if 'r' is in pos 3 and 5, but we know there is only one 'r'. So "} \\
            }
        }
    }
    & & \\ \\

    \theutterance \stepcounter{utterance}  
    & & & \multicolumn{2}{p{0.3\linewidth}}{
        \cellcolor[rgb]{0.9,0.9,0.9}{
            \makecell[{{p{\linewidth}}}]{
                \texttt{\tiny{[GM$|$GM]}}
                \texttt{Error: The response should contain the 'guess:' keyword exactly once.} \\
            }
        }
    }
    & & \\ \\

    \theutterance \stepcounter{utterance}  
    & & & \multicolumn{2}{p{0.3\linewidth}}{
        \cellcolor[rgb]{0.9,0.9,0.9}{
            \makecell[{{p{\linewidth}}}]{
                \texttt{\tiny{[GM$|$GM]}}
                \texttt{game\_result = ABORT} \\
            }
        }
    }
    & & \\ \\

    \theutterance \stepcounter{utterance}  
    & & \multicolumn{4}{p{0.6\linewidth}}{
        \cellcolor[rgb]{0.95,0.95,0.95}{
            \makecell[{{p{\linewidth}}}]{
                \texttt{\tiny{[P1$\leftrightarrow$P]}}
                \texttt{Post{-}episode Reflexion (not sent to GM):} \\
\texttt{New plan: The previous attempt failed because I generated a response that contained internal monologue, hesitation ("Wait", "Let's try"), and incomplete sentences, violating the strict "explanation: details \textbackslash{}n guess: word" format. The Game Master rejected the message before it could be evaluated as a guess. Next time, I will strictly separate my reasoning from the output. I will perform all logical deductions (checking letter constraints, verifying word validity against red letters, and ensuring the word is 5 letters) internally before generating the response. The final output will contain only a concise explanation of the logic used to select the word and the word itself, with no conversational filler, self{-}correction, or trailing text. Specifically, I will ensure the explanation ends with a period and the guess is a single, valid 5{-}letter word on the next line.} \\
            }
        }
    }
    & \\ \\

\end{supertabular}
}

\end{document}
